\section{Instrument gate, pilot, confirmatory study, and stopping rules}
\label{sec:stages}

\paragraph{The instrument gate comes first, and it is the study's central risk.}
Because the design is exhaustive re-execution rather than randomized treatment
(\cref{sec:causal}), every estimand and \cref{thm:safe} itself depend on
branches beginning from an identical restored state.  Runtime integration
therefore starts with one minimal toy task, one restored checkpoint, and three
isolated branches, and must demonstrate restoration of the five registered
layers --- workspace, executable environment, session history, budget counters,
and declared external tool state --- with the feasible mask reproduced under its
own check rather than as a sixth layer; isolation of branches
against cross-branch writes \emph{and} against shared external state such as
local inference caches, shared process state, and shared tool sessions; evaluator
integrity; and trace completeness.  This stage validates infrastructure only and
supplies no scientific evidence.  Its result is reported whichever way it falls:
a failed gate is a finding about benchmark readiness, not a reason to weaken
\cref{sec:causal}.

If the gate passes, a later engineering dry run may validate intervention
semantics and cost capture; it cannot estimate effects or rank policies.  A
developmental pilot may then estimate eligible-checkpoint yield, restoration and
instrumentation reliability, availability per regime, raw cost distributions,
adjudicator agreement, paired root-level variance, within-root stochastic
variance, and whether recovery semantics remain coherent.  Pilot data are
excluded from confirmatory inference.  For E3, the approved claim and root-count
rule are frozen before the pilot and its output is only a go/no-go assessment;
failure requires a new scientific approval rather than threshold adaptation.
For E4, \cref{eq:sample-ceiling} separately fixes the certification ceiling and
the pilot reports whether the achievable root count clears it.

Before pair construction, a seeded group-disjoint allocation carves a
\emph{pair-development} subset from the training acquisition pool; the
remainder is the model-training pool.  Only pair-development data may
instantiate $P$, $\lambda_0$, and $\lambda_P$, and the split manifest and exact
pair are hash-bound before the pilot.  Pilot, selection, calibration, and final
roots remain disjoint from both training subsets.

After the pilot, the finite candidate class and the model-training, selection,
calibration, and untouched final-evaluation splits are frozen.  Candidate
training uses only the model-training remainder; selection and simultaneous
certification use their declared held-out splits.  Reporting a pre-registered comparator on the
final split is not selecting on it: only choices made after seeing final data
trigger simultaneous coverage.  The confirmatory estimand \eqref{eq:personalization}
is one pre-registered contrast, and the final split contributes a single
ordinary confidence interval for it --- a curve over the registered prevalence
grid, with the guarantee interval of \cref{cor:interval} marked.  The remaining
members of the comparator hierarchy of \cref{sec:policy} are reported on the
final split descriptively, with point estimates and, where prespecified, their
own intervals, but no selection among them is made from those data; if a
comparison or prevalence point not registered in advance is added on final data,
simultaneous coverage is constructed for that selection instead.  Root and
template clustering, unequal availability, missing checkpoint blocks, and all
certification conditions enter the confirmatory sample projection.

Training the parametric policies of $\Dcal$ is permitted entirely inside the
model-training remainder, as \cref{sec:rl} specifies; it does not reintroduce sequential
positivity, because the decision remains one-step and every feasible arm is
branched.  What the frozen protocol prohibits is precise: re-fitting any
candidate, adjusting the registered gate, relaxing a constraint, or reselecting
the prevalence point after the candidate class is frozen.  The gate is a
registered procedure --- applying it on calibration data produces the derived
deployment rules $d^\dagger$ and $d_{\mathrm{C,cert}}$
of \cref{sec:policy}, which is the certification step, not adaptive reuse.  The
model-training remainder fits the unconstrained and constrained policy members; the selection split fixes architecture,
hyperparameters, checkpoint and the optimistic controller's point estimates; the
calibration split constructs the simultaneous bounds, forms the gate of
\cref{cor:proposal}, and determines the certified set; the untouched final split
evaluates the frozen policy alone, after which no policy swap, gate adjustment,
constraint relaxation, or prevalence reselection is permitted.

The branching data would be a full-feedback contextual-bandit benchmark ---
every feasible action observed conditional on pre-treatment feasibility, with
clean counterfactuals and randomised fault assignment --- and their release is a
conditional deliverable of the program rather than a claimed contribution of
this paper.  Release is nevertheless conditional on licensing, runtime and model
terms, storage and versioning, and contamination control; if any of those
conditions cannot be met, the paper reports which one blocked release rather
than silently dropping the deliverable.

\paragraph{Claim hierarchy.}
In the retained historical plan E3 was primary, E4 the safety claim, E1/E2
supporting, and E5 robustness.  All of these are now deferred; the completed
claims of this report are M1--M3, and none of the studies below has run.  E3 and
E4 retain separate feasibility conditions and verdicts.  E4 is deferred from
the execution program by decision, not by blocker; its theory is
retained unchanged and its conditions below are stated for a reinstatement.
E3 is deferred by the same route decision and is not the current primary claim.

\emph{Ordinal sufficiency} tests the null \eqref{eq:null} by budget-matched
capacity-$q$ selection under a risk score against a cross-fitted $a_v$, on the
domain where $g_0$ and $g_v$ are jointly available, across a declared capacity
grid.  It needs no recovery regime, so it stays identified even if recovery
semantics prove incoherent, and its null is sharp rather than conventional; it
is the cheapest identified comparison in the design and not a fallback.

\emph{The E4 safety claim, certified personalization,} tests whether
\eqref{eq:personalization} is positive on the final holdout while the constraints of \cref{prop:certification}
are met, with the guarantee interval of \cref{cor:interval} marked on the
reported prevalence curve.  Its endpoints include the fallback rate of
\eqref{eq:gated-rule} --- a controller that never deviates is safe and useless
--- and the no-fault harm fraction \eqref{eq:harm-fraction}, a primary endpoint
and not a robustness check.

\paragraph{Numeric gates.}
The full design proceeds only if all of the following hold.  These values are
registered in advance and are not to be relaxed afterwards.  G1 is measured on
a dedicated pre-pilot instrument pool, not on pilot data.

\begin{enumerate}[label=(G\arabic*),leftmargin=3em]
\item \textbf{Restoration.} At least $0.95$ of attempted checkpoint restoration
  \emph{blocks} --- a block being the parent together with every fresh sibling ---
  reproduce all five registered layers exactly: workspace, executable
  environment, session history, budget counters, and declared external tool
  state.  The availability mask is not a layer and is reproduced under its own
  check.  Zero detected cross-branch writes and zero detected shared-state leaks
  are required across at least $200$ valid completed branch triples, with
  complete passing isolation probes and no unresolved required check.  Absence
  of evidence is never zero leakage.
\item \textbf{Availability.} Each regime class is available at at least $0.30$ of
  eligible checkpoints, and the three-way common-availability domain covers at
  least $0.20$, so that action-ranking contrasts are estimable rather than
  nominal.
\item \textbf{E4 certification ceiling} \emph{(if E4 is reinstated).} The projected E4 confirmatory root count clears the
  value \eqref{eq:sample-ceiling} returns at the registered class sizes and
  margins $(|\Dcal|,B,K,\delta,\eta)$, after inflation for root and template
  clustering and for invalidated checkpoint blocks.
\item \textbf{Adjudication.} Independent adjudicator agreement on terminal
  category is at least $\kappa=0.75$, with execution correctness and scientific
  validity scored separately.
\item \textbf{Recovery coherence.} For each fault family, the recovery
  precondition, rollback boundary, recomputation rule, and machine-checkable
  postcondition are demonstrated on at least $20$ checkpoints, with a
  postcondition false-pass rate below $0.05$.
\end{enumerate}

\paragraph{Stopping and reframing.}
If (G1) fails, no confirmatory causal claim is made and the paper reports the
instrument result.  If (G2), (G3), or (G5) fails while $g_0$ and $g_v$
remain available and the approved E3 precision target remains feasible, the
study may proceed with ordinal sufficiency alone and the certified-personalization
claim is withdrawn, not weakened.  If the
personalization estimand \eqref{eq:personalization} is non-positive on the final
holdout, or the no-fault constraint cannot be certified, the successful-policy
claim is stopped and the work is reported as causal characterization.

The stopping rule keeps its force, and what it gates is the framing, not the
paper.  If the null \eqref{eq:null} is rejected, the actionability gap exists
and risk-ranked budgets misallocate.  If it is not rejected, the study may say
only that ordinal insufficiency was not established at the registered
capacities and precision; it reports simultaneous confidence intervals and the
design's power against prespecified alternatives.  Failure to reject cannot
establish ordinal sufficiency.  An approximate-sufficiency claim would require
a separately registered equivalence tolerance and simultaneous upper confidence
bounds below that tolerance across the capacity grid.  Neither outcome is a
failure of the study, but only rejection licenses the insufficiency
claim; the E0-M, E2 and harm measurements do not depend on it.  With E3
deferred, this stopping rule governs a future execution and does not gate the
completed measurement panels.
